\section*{Глава 3. Разработка веб-приложения}
\addcontentsline{toc}{section}{Глава 3. Разработка веб-приложения}
\stepcounter{section}


\subsection{Организация структуры проекта и контейнеризация}

Проект организован в виде монорепозитория, содержащего три
независимых компонента: серверное API на Go (\texttt{plantation-api}),
клиентское приложение на React (\texttt{plantation-front}) и
конфигурацию оркестрации Docker~Compose. Такая структура
позволяет вести параллельную разработку серверной и клиентской
частей, сохраняя единую точку развёртывания.

\begin{lstlisting}[caption={Файловая структура проекта}, label={lst:project_tree}]
plantation/
  docker-compose.yml
  plantation-api/
    cmd/api/main.go
    internal/
      handler/     (auth, sector, plant, watering, training, files, ...)
      middleware/   (auth.go)
      model/        (models.go)
      simulation/   (engine.go, formulas.go)
      storage/      (postgres.go, seed.go)
      ws/           (hub.go)
      xlsx/         (xlsx.go)
    migrations/     (001_init.sql, 002_simulation.sql,
                     003_gamification.sql)
    Dockerfile
    go.mod
  plantation-front/
    src/
      pages/        (Dashboard, Sectors, Watering, Plants, ...)
      components/   (Shell, visuals, icons)
      api.js
      App.jsx
      main.jsx
    package.json
    index.html
\end{lstlisting}

Развёртывание реализовано с использованием
Docker~Compose~\cite{docker_compose_docs,docker_compose_file_reference},
объединяющего три сервиса. Сервис \texttt{db} использует образ
\texttt{postgres:16-alpine}, автоматически исполняя SQL-миграции
из каталога \texttt{migrations} при первом запуске через
механизм \texttt{docker-entrypoint-initdb.d}.
Сервис \texttt{api} собирается из многоступенчатого
Dockerfile~\cite{docker_multistage_docs} (листинг~\ref{lst:dockerfile}),
что позволяет отделить среду компиляции (образ \texttt{golang:1.23-alpine})
от финального образа (\texttt{alpine:3.20}), сокращая размер
контейнера до $\approx$\,20~МБ. Сервис \texttt{front}
запускает dev-сервер Vite~\cite{vite_official}, проксирующий
запросы \texttt{/api} и \texttt{/ws} к сервису \texttt{api}.

\begin{lstlisting}[caption={Многоступенчатая сборка API-сервера (Dockerfile)}, label={lst:dockerfile}]
FROM golang:1.23-alpine AS builder
WORKDIR /app
COPY go.mod go.sum ./
RUN go mod download
COPY . .
RUN CGO_ENABLED=0 go build -o /bin/api ./cmd/api

FROM alpine:3.20
RUN apk add --no-cache ca-certificates
COPY --from=builder /bin/api /bin/api
EXPOSE 8080
CMD ["/bin/api"]
\end{lstlisting}

Зависимости между сервисами задаются директивой
\texttt{depends\_on} с условием \texttt{service\_healthy}:
API-сервер стартует только после прохождения healthcheck
базы данных (\texttt{pg\_isready}), что исключает сбои
при подключении к СУБД.


\subsection{Реализация базы данных и миграций}

Структура базы данных реализована посредством трёх SQL-миграций,
исполняемых PostgreSQL~16~\cite{postgresql_official,postgresql_16_release_notes}
при инициализации контейнера. Первая миграция
(\texttt{001\_init.sql}) создаёт основные таблицы: пользователей
с перечислимым типом роли (\texttt{user\_role}), секторов
плантации, растений и журнала полива. Вторая миграция
(\texttt{002\_simulation.sql}) добавляет таблицы телеметрии
и уведомлений, а также расширяет таблицу секторов полями
для суточного лимита воды. Третья миграция
(\texttt{003\_gamification.sql}) добавляет таблицы обучающей
аналитики (\texttt{training\_scores}) и профилей генератора
погоды (\texttt{weather\_configs}), а также расширяет таблицу
секторов полями индекса водного стресса \texttt{cwsi},
счётчиков серий для начисления бейджей (\texttt{healthy\_streak},
\texttt{safe\_streak}, \texttt{crisis\_streak}), блокировки
полива при поломке оборудования
(\texttt{equipment\_locked\_ticks}), флага активного нашествия
вредителей (\texttt{pest\_active}) и типа последнего поднятого
уведомления (\texttt{last\_alert\_kind}). Та же схема идемпотентно
применяется при старте сервера методом
\texttt{storage.EnsureSchema} (через
\texttt{ADD COLUMN IF NOT EXISTS} и
\texttt{CREATE TABLE IF NOT EXISTS}), что обеспечивает
совместимость с уже существующими базами данных без потери
данных.

\begin{lstlisting}[caption={Создание таблицы секторов (фрагмент 001\_init.sql)}, label={lst:sql_sectors}]
CREATE TABLE sectors (
  id              UUID PRIMARY KEY DEFAULT gen_random_uuid(),
  name            VARCHAR(255) NOT NULL,
  area_sqm        FLOAT NOT NULL DEFAULT 0,
  soil_moisture   FLOAT NOT NULL DEFAULT 50,
  temperature     FLOAT NOT NULL DEFAULT 25,
  health_index    FLOAT NOT NULL DEFAULT 1.0,
  gdd_cumulative  FLOAT NOT NULL DEFAULT 0,
  phenophase      VARCHAR(10) NOT NULL DEFAULT '00',
  ks_water        FLOAT NOT NULL DEFAULT 1.0,
  ks_aeration     FLOAT NOT NULL DEFAULT 1.0,
  deficit_dr      FLOAT NOT NULL DEFAULT 0,
  status          VARCHAR(50) NOT NULL DEFAULT 'normal',
  operator_id     UUID REFERENCES users(id) ON DELETE SET NULL,
  created_at      TIMESTAMPTZ NOT NULL DEFAULT now(),
  updated_at      TIMESTAMPTZ NOT NULL DEFAULT now()
);
\end{lstlisting}

Таблица телеметрии индексирована составным индексом
(\texttt{sector\_id}, \texttt{recorded\_at~DESC}), что
обеспечивает эффективную выборку временных рядов для
построения графиков в интерфейсе.

Взаимодействие с базой данных реализовано через ORM-библиотеку
Bun~\cite{bun_uptrace_official,bun_github}. Модели данных
описаны в пакете \texttt{internal/model} с использованием
тегов \texttt{bun:} для маппинга полей Go-структур на столбцы
PostgreSQL. Пакет \texttt{internal/storage} содержит методы
хранилища с типобезопасными fluent-запросами
(листинг~\ref{lst:storage_sample}).

\begin{lstlisting}[caption={Пример метода хранилища (выборка телеметрии)}, label={lst:storage_sample}]
func (s *Storage) ListTelemetry(
    ctx context.Context, sectorID string, limit int,
) ([]model.Telemetry, error) {
    var rows []model.Telemetry
    q := s.DB.NewSelect().
        Model(&rows).
        Order("recorded_at DESC").
        Limit(limit)
    if sectorID != "" {
        q = q.Where("sector_id = ?", sectorID)
    }
    err := q.Scan(ctx)
    return rows, err
}
\end{lstlisting}


\subsection{Реализация серверного API и маршрутизации}

Серверная часть реализована на языке Go~1.23~\cite{go_dev_official,effective_go}
с использованием маршрутизатора go-chi~\cite{go_chi_github,go_chi_pkgdev}.
Точка входа (\texttt{cmd/api/main.go}) выполняет инициализацию
хранилища, создание WebSocket-хаба, запуск движка симуляции
и конфигурацию HTTP-маршрутов.

Маршруты разделены на три группы:
\begin{enumerate}
    \item \textbf{Публичные} --- регистрация и аутентификация
    (\texttt{POST /api/auth/register}, \texttt{POST /api/auth/login}).
    \item \textbf{Защищённые} (требуют JWT-токен) --- чтение
    секторов, растений, телеметрии, уведомлений и отчётов.
    \item \textbf{Административные} (только для роли
    \texttt{agronomist}) --- создание, изменение, удаление и
    ручная корректировка секторов, назначение и снятие
    операторов, импорт/экспорт CSV, экспорт телеметрии в XLSX
    и настройка параметров генератора погоды.
\end{enumerate}

\begin{lstlisting}[caption={Конфигурация защищённых маршрутов (main.go)}, label={lst:routes}]
r.Group(func(r chi.Router) {
    r.Use(middleware.Auth(secret))

    r.Get("/api/sectors", sectorH.List)
    r.Get("/api/sectors/{id}", sectorH.Get)
    r.Get("/api/sectors/my", sectorH.ListMy)

    // agronomist only
    r.Group(func(r chi.Router) {
        r.Use(middleware.RequireRole("agronomist"))
        r.Post("/api/sectors", sectorH.Create)
        r.Put("/api/sectors/{id}", sectorH.Update)
        r.Delete("/api/sectors/{id}", sectorH.Delete)
        r.Patch("/api/sectors/{id}/override", sectorH.Override)
        r.Put("/api/sectors/{id}/assign", sectorH.Assign)
        r.Delete("/api/sectors/{id}/assign", sectorH.Unassign)
        r.Get("/api/export/sectors", fileH.ExportSectors)
        r.Get("/api/export/telemetry/{id}", fileH.ExportTelemetry)
        r.Post("/api/import/sectors", fileH.ImportSectors)
        r.Get("/api/weather-config", trainingH.GetWeatherConfig)
        r.Put("/api/weather-config", trainingH.UpdateWeatherConfig)
        r.Post("/api/training/{userId}/award", trainingH.Award)
    })

    r.Post("/api/water", waterH.Water)
    r.Post("/api/sectors/{id}/treat", waterH.Treat)
    r.Get("/api/telemetry/{sectorId}", reportH.Telemetry)
    r.Get("/api/reports/{sectorId}", reportH.Summary)
    r.Get("/api/notifications", notifH.List)
    r.Get("/api/leaderboard", trainingH.Leaderboard)
    r.Get("/api/score/my", trainingH.MyScore)
})
\end{lstlisting}

В таблице~\ref{tab:api_endpoints} приведён полный перечень
API-эндпоинтов с указанием HTTP-метода, пути, требуемой роли
и назначения.

\begin{longtable}{|l|l|l|p{4.8cm}|}
    \caption{Перечень REST API-эндпоинтов} \label{tab:api_endpoints} \\

    \hline
    \textbf{Метод} & \textbf{Путь} & \textbf{Роль} & \textbf{Назначение} \\ \hline
    \endfirsthead

    \tabcont{\ref{tab:api_endpoints}}
    \textbf{Метод} & \textbf{Путь} & \textbf{Роль} & \textbf{Назначение} \\ \hline
    \endhead

    \hline
    \endfoot

    POST & /api/auth/register & --- & Регистрация пользователя \\ \hline
    POST & /api/auth/login    & --- & Аутентификация, выдача JWT \\ \hline
    GET  & /api/sectors       & любая & Список секторов \\ \hline
    GET  & /api/sectors/\{id\}  & любая & Детали сектора \\ \hline
    GET  & /api/sectors/my    & любая & Секторы текущего оператора \\ \hline
    POST & /api/sectors       & agro  & Создание сектора \\ \hline
    PUT  & /api/sectors/\{id\}  & agro  & Обновление сектора \\ \hline
    DELETE & /api/sectors/\{id\} & agro & Удаление сектора \\ \hline
    PUT  & /api/sectors/\{id\}/assign & agro & Назначение оператора \\ \hline
    DELETE & /api/sectors/\{id\}/assign & agro & Снятие оператора \\ \hline
    PATCH & /api/sectors/\{id\}/override & agro & Ручная корректировка параметров и инъекция событий \\ \hline
    GET  & /api/plants        & любая & Список растений \\ \hline
    POST & /api/plants        & любая & Добавление растения \\ \hline
    DELETE & /api/plants/\{id\} & любая & Удаление растения \\ \hline
    POST & /api/water         & любая & Выполнение полива \\ \hline
    GET  & /api/water/stats/\{id\} & любая & Статистика полива \\ \hline
    POST & /api/sectors/\{id\}/treat & любая & Обработка сектора от вредителей \\ \hline
    GET  & /api/notifications & любая & Список уведомлений \\ \hline
    PUT  & /api/notifications/\{id\}/read & любая & Пометить как прочитано \\ \hline
    GET  & /api/telemetry/\{id\} & любая & Временной ряд телеметрии \\ \hline
    GET  & /api/reports/\{id\} & любая & Агрегированный отчёт \\ \hline
    GET  & /api/export/sectors & agro & Экспорт секторов в CSV \\ \hline
    GET  & /api/export/telemetry/\{id\} & agro & Экспорт телеметрии сектора в XLSX \\ \hline
    POST & /api/import/sectors & agro & Импорт секторов из CSV \\ \hline
    GET  & /api/weather-config & agro & Параметры генератора погоды \\ \hline
    PUT  & /api/weather-config & agro & Настройка генератора погоды и вероятностей событий \\ \hline
    GET  & /api/leaderboard & любая & Таблица лидеров операторов \\ \hline
    GET  & /api/score/my & любая & Личный счёт и бейджи оператора \\ \hline
    POST & /api/training/\{id\}/award & agro & Ручное начисление баллов и бейджей оператору \\ \hline
\end{longtable}

CORS-политика настроена через middleware \texttt{go-chi/cors}
и разрешает запросы с локальных адресов разработки,
передачу заголовков \texttt{Authorization} и
\texttt{Content-Type}, а также использование методов
\texttt{GET}, \texttt{POST}, \texttt{PUT}, \texttt{PATCH},
\texttt{DELETE} и \texttt{OPTIONS}.


\subsection{Реализация аутентификации и ролевой модели доступа}

Модуль аутентификации реализован в соответствии с проектными
решениями, описанными в подразделе~1.2.5. Пароли хешируются
алгоритмом bcrypt~\cite{provos_mazieres_1999,pkggo_bcrypt} со
стандартным коэффициентом стоимости (10 раундов). При успешной
аутентификации сервер генерирует JWT-токен
(RFC~7519)~\cite{rfc7519_ietf,jwt_io_intro} с подписью
HMAC-SHA256, содержащий поля \texttt{sub} (UUID пользователя),
\texttt{role} (\texttt{agronomist} или \texttt{operator}) и
\texttt{exp} (срок действия --- 24~часа).

Middleware-функция \texttt{Auth} (листинг~\ref{lst:auth_mw})
перехватывает каждый запрос к защищённым маршрутам, извлекает
токен из заголовка \texttt{Authorization: Bearer~<token>},
верифицирует подпись и внедряет идентификатор и роль
пользователя в контекст запроса Go~\cite{go_context_pkgdev}.

\begin{lstlisting}[caption={Middleware аутентификации (middleware/auth.go)}, label={lst:auth_mw}]
func Auth(secret string) func(http.Handler) http.Handler {
    return func(next http.Handler) http.Handler {
        return http.HandlerFunc(func(w http.ResponseWriter,
            r *http.Request) {
            header := r.Header.Get("Authorization")
            if !strings.HasPrefix(header, "Bearer ") {
                http.Error(w,
                    `{"error":"unauthorized"}`,
                    http.StatusUnauthorized)
                return
            }
            tokenStr := strings.TrimPrefix(header, "Bearer ")
            token, err := jwt.Parse(tokenStr,
                func(t *jwt.Token) (any, error) {
                    return []byte(secret), nil
                })
            if err != nil || !token.Valid {
                http.Error(w,
                    `{"error":"invalid token"}`,
                    http.StatusUnauthorized)
                return
            }
            claims := token.Claims.(jwt.MapClaims)
            ctx := context.WithValue(r.Context(),
                UserIDKey, claims["sub"].(string))
            ctx = context.WithValue(ctx,
                RoleKey, claims["role"].(string))
            next.ServeHTTP(w, r.WithContext(ctx))
        })
    }
}
\end{lstlisting}

Дополнительный middleware \texttt{RequireRole} принимает
список допустимых ролей и возвращает HTTP~403 (Forbidden),
если роль текущего пользователя не входит в разрешённый
набор. Это обеспечивает защиту административных маршрутов
(создание и удаление секторов, назначение операторов,
импорт/экспорт) на уровне серверного API.


\subsection{Реализация движка имитационной модели}

Движок симуляции реализован в пакете \texttt{internal/simulation}
и представляет собой тиковую архитектуру, описанную в
подразделе~2.4. При старте сервера создаётся экземпляр
\texttt{Engine}, запускающий горутину~\cite{effective_go} с
периодическим таймером (\texttt{time.Ticker}) с интервалом
10~секунд. На каждом тике движок сначала генерирует единые
для всей плантации погодные условия (один климат на
плантацию), затем извлекает из базы данных все секторы и
обрабатывает их \textbf{параллельно} --- каждый сектор в
отдельной горутине с синхронизацией через \texttt{sync.WaitGroup},
что отражает независимость секторов как сущностей. Начисление
очков операторам применяется последовательно после барьера
синхронизации во избежание состояний гонки при записи
(корректность подтверждена детектором гонок \texttt{go~-race}).
Для воспроизводимости обучающих сценариев предусмотрен
детерминированный режим, активируемый переменной окружения
\texttt{SIM\_SEED}.

\subsubsection{Конвейер подсистем}

Тик движка реализует полный конвейер из
девяти подсистем, описанных на рисунке~\ref{fig:tick_pipeline}:

\begin{enumerate}
    \item \textbf{WeatherTick} --- генерация единых для всей
    плантации погодных условий. Состояние дня
    (<<сухой/влажный>>) определяется цепью Маркова первого
    порядка (формула~(\ref{eq:markov_weather})) с матрицей
    переходов $P_{01} = 0{,}20$ (сухой $\to$ влажный) и
    $P_{11} = 0{,}55$ (влажный $\to$ влажный); вероятности
    настраиваются агрономом через профиль
    \texttt{weather\_configs}. Суточная температура
    моделируется авторегрессионным процессом первого порядка
    AR(1), условным на состояние дня (во влажные дни прохладнее
    и меньше суточная амплитуда). Осадки во влажный день ---
    выборка из гамма-распределения $\Gamma(\alpha, \beta)$
    методом Марсальи--Цанга. В сухой день с вероятностью
    $\approx$\,5\% методом Монте-Карло генерируется аномальная
    жара, повышающая дневной максимум температуры.
    \item \textbf{ET0Tick} --- расчёт эталонной эвапотранспирации
    $ET_0$. Реализованы два метода, переключаемые полем
    \texttt{et\_method}: упрощённое уравнение Харгривса--Самани
    (по умолчанию, формула~(\ref{eq:hargreaves})) и полное
    уравнение Пенмана--Монтейта (FAO-56). Внеатмосферная
    радиация $R_a$ вычисляется по широте и порядковому номеру
    дня в году, а не задаётся константой.
    \item \textbf{WaterBalanceTick} --- обновление дефицита
    влаги по уравнению водного баланса
    (формула~(\ref{eq:water_balance})). Объём полива $I(t)$
    применяется непосредственно обработчиком API при действии
    оператора; движок обрабатывает естественные процессы:
    эффективные осадки $P_{eff}$, эвапотранспирацию $ET_c$
    и глубинную перколяцию $DP$.
    \item \textbf{PhenologyTick} --- накопление суммы
    эффективных температур GDD
    (формула~(\ref{eq:gdd})) и переключение фенофазы
    по расширенной шкале BBCH~\cite{wei2013_litchi_bbch}.
    Прирост GDD масштабируется тем же коэффициентом сжатия
    времени, что эвапотранспирация и осадки, поэтому фенофазы
    сменяются постепенно, а не мгновенно.
    \item \textbf{StressTick} --- расчёт коэффициентов
    стресса: водного $K_s$ (формула~(\ref{eq:ks_stress})),
    аэрационного $K_{s,aer}$
    (формула~(\ref{eq:ks_aer})) и теплового $K_{s,T}$.
    Дополнительно вычисляется индекс водного стресса культуры
    CWSI (по базовым линиям Идсо--Джексона) для приборной
    панели агронома.
    \item \textbf{HealthTick} --- обновление индекса здоровья
    $H$ по формуле~(\ref{eq:health_index}).
    \item \textbf{EventTick} --- стохастические события
    методом Монте-Карло
    (формула~(\ref{eq:event_probability})). \emph{Нашествие
    вредителей} (вероятность растёт с накоплением GDD)
    переводит сектор в состояние стойкого заражения
    (\texttt{pest\_active}), которое наносит постоянный урон
    индексу здоровья ($\alpha_{pest} = 0{,}03$ на тик) и
    блокирует восстановление до тех пор, пока оператор не
    выполнит \emph{обработку} (\texttt{POST /api/sectors/\{id\}/treat});
    при наступлении формируется уведомление \texttt{pest\_attack}.
    \emph{Поломка оборудования} не штрафует здоровье напрямую,
    а блокирует полив сектора на 1--3 тика (поле
    \texttt{equipment\_locked\_ticks}), вынуждая оператора
    планировать действия. Вероятности всех событий настраиваются
    агрономом через профиль \texttt{weather\_configs}.
    \item \textbf{StatusTick} --- определение статуса
    сектора по автомату состояний
    (рисунок~\ref{fig:state_machine}).
    \item \textbf{ScoringTick} --- начисление очков и
    присуждение бейджей оператору, ответственному за сектор
    (подробнее --- подраздел~<<Реализация геймифицированного
    интерфейса>>). Применяется последовательно после
    параллельной обработки всех секторов.
\end{enumerate}

\subsubsection{Реализация фенологических переходов}

Фенофазы описаны таблицей пороговых значений кумулятивных
GDD с соответствующим базальным коэффициентом культуры
$K_{cb}$ (листинг~\ref{lst:pheno_table}). При достижении
порогового значения GDD код фенофазы автоматически
обновляется, а коэффициент $K_{cb}$ используется в формуле
дуального коэффициента культуры
(формула~(\ref{eq:dual_kc})):

\begin{equation}
    ET_c = (K_s \cdot K_{cb}(\text{фенофаза}) + K_e) \cdot ET_0.
    \label{eq:etc_impl}
\end{equation}

\begin{lstlisting}[caption={Таблица фенологических порогов (engine.go)}, label={lst:pheno_table}]
var phenoPhases = []struct {
    minGDD float64
    code   string
    kcb    float64
}{
    {0, "00", 0.45},     // dormancy
    {150, "01", 0.55},   // bud break
    {400, "03", 0.70},   // shoot development
    {700, "05", 0.85},   // inflorescence emergence
    {1000, "06", 0.95},  // flowering
    {1500, "07", 1.00},  // fruit development
    {2200, "08", 0.85},  // ripening
}
\end{lstlisting}

\subsubsection{Реализация расчёта индекса здоровья}

Индекс здоровья $H$ обновляется на каждом тике по
формуле~(\ref{eq:health_index}). В
листинге~\ref{lst:health_calc} представлен фрагмент
реализации, в котором коэффициенты деградации
$\alpha_{drought} = 0{,}03$, $\alpha_{flood} = 0{,}02$,
$\alpha_{heat} = 0{,}01$, $\alpha_{pest} = 0{,}03$ (урон от
активного нашествия вредителей) и коэффициент восстановления
$\beta_{rec} = 0{,}05$ заданы как константы пакета.
Восстановление активируется только при отсутствии всех
видов стресса ($K_s > 0{,}9$, $K_{s,aer} > 0{,}9$,
$K_{s,T} > 0{,}9$) и отсутствии заражения вредителями.

\begin{lstlisting}[caption={Расчёт индекса здоровья (engine.go)}, label={lst:health_calc}]
health := sec.HealthIndex
deltaDrought := alphaDrought * (1.0 - sec.KsWater)
deltaFlood   := alphaFlood * (1.0 - sec.KsAeration)
deltaHeat    := alphaHeat * (1.0 - ksT)
deltaPest    := 0.0
if sec.PestActive {
    deltaPest = alphaPest
}

rec := 0.0
if sec.KsWater > 0.9 && sec.KsAeration > 0.9 &&
    ksT > 0.9 && !sec.PestActive {
    rec = betaRec * sec.KsWater *
        sec.KsAeration * (1.0 - health)
}

health = health - deltaDrought - deltaFlood -
    deltaHeat - deltaPest + rec
health = math.Max(0, math.Min(1.0, health))
sec.HealthIndex = health
\end{lstlisting}

\subsubsection{Реализация автомата состояний}

Статус сектора определяется по приоритетной цепочке
условий, соответствующей диаграмме состояний UML
(рисунок~\ref{fig:state_machine}). Реализованы восемь
состояний: \texttt{normal}, \texttt{drought},
\texttt{overwatered}, \texttt{heat\_stress},
\texttt{pest} (активное нашествие вредителей),
\texttt{critical}, \texttt{recovering} и \texttt{dead}.
Переход в состояние \texttt{recovering} активируется
при восстановлении стрессовых коэффициентов до уровня
$> 0{,}85$ при индексе здоровья в диапазоне [0{,}3; 0{,}7).
Состояние \texttt{dead} фиксируется при падении индекса
здоровья ниже 0{,}1 и генерирует уведомление о необходимости
пересадки.


\subsection{Реализация обмена данными в реальном времени}

Полнодуплексная связь между сервером и клиентами реализована
по протоколу WebSocket (RFC~6455)~\cite{rfc6455_ietf,mdn_websocket}
с использованием библиотеки
gorilla/websocket~\cite{gorilla_websocket_github,gorilla_websocket_pkgdev}.

\subsubsection{Серверный WebSocket-хаб}

Структура \texttt{Hub} (листинг~\ref{lst:ws_hub}) реализует
паттерн <<publish-subscribe>>. Поскольку движок симуляции
рассылает обновления из множества горутин секторов
одновременно, а библиотека \texttt{gorilla/websocket}
запрещает конкурентную запись в одно соединение,
потокобезопасность обеспечивается двумя мьютексами: карта
клиентов защищена \texttt{sync.RWMutex}, а фактические записи
в сокеты сериализуются отдельным мьютексом
\texttt{writeMu}. Рассылка сначала снимает <<снимок>> набора
клиентов под разделяемой блокировкой (не удерживая её во время
сетевого ввода-вывода), после чего последовательно
отправляет сообщение каждому соединению. При подключении
нового клиента горутина \texttt{ReadPump} поддерживает
соединение живым, а при разрыве --- удаляет клиента из карты.

\begin{lstlisting}[caption={WebSocket-хаб: широковещательная рассылка (ws/hub.go)}, label={lst:ws_hub}]
func (h *Hub) Broadcast(event string, data any) {
    msg, err := json.Marshal(map[string]any{
        "event": event,
        "data":  data,
    })
    if err != nil {
        return
    }
    // snapshot clients without holding the lock during I/O
    h.mu.RLock()
    conns := make([]*websocket.Conn, 0, len(h.clients))
    for conn := range h.clients {
        conns = append(conns, conn)
    }
    h.mu.RUnlock()

    // serialise writes: gorilla forbids concurrent
    // writes to a single connection
    h.writeMu.Lock()
    defer h.writeMu.Unlock()
    for _, conn := range conns {
        if err := conn.WriteMessage(
            websocket.TextMessage, msg); err != nil {
            conn.Close()
            h.mu.Lock()
            delete(h.clients, conn)
            h.mu.Unlock()
        }
    }
}
\end{lstlisting}

Движок симуляции вызывает \texttt{Hub.Broadcast} с двумя
типами событий:
\begin{itemize}
    \item \texttt{sector:update} --- обновлённое состояние
    сектора после каждого тика;
    \item \texttt{notification} --- критическое уведомление
    (засуха, переувлажнение, падение здоровья, нашествие
    вредителей, поломка оборудования).
\end{itemize}

Во избежание <<спама>> уведомления по водно-стрессовым
состояниям формируются только при \emph{смене} состояния
сектора, а не на каждом тике: тип последнего поднятого
уведомления хранится в поле \texttt{last\_alert\_kind} и
сравнивается с текущим. Обработчик полива дополнительно
рассылает событие \texttt{sector:watered} с информацией о
выполненном поливе.

\subsubsection{Клиентское WebSocket-подключение}

На стороне клиента функция \texttt{connectWS}
(листинг~\ref{lst:ws_client}) устанавливает WebSocket-соединение
с автоматическим переподключением через 3~секунды при обрыве.
Протокол определяется динамически (\texttt{ws:} или
\texttt{wss:}) в зависимости от протокола загрузки страницы.

\begin{lstlisting}[caption={Клиентское WebSocket-подключение (api.js)}, label={lst:ws_client}]
export function connectWS(onMessage, onStatus) {
    let ws = null
    let reconnectTimer = null
    let closed = false
    function connect() {
        const proto = window.location.protocol ===
            'https:' ? 'wss:' : 'ws:'
        ws = new WebSocket(
            `${proto}//${window.location.host}/ws`)
        ws.onopen = () => { onStatus?.(true) }
        ws.onmessage = e => {
            try { onMessage(JSON.parse(e.data)) }
            catch {}
        }
        ws.onclose = () => {
            onStatus?.(false)
            if (!closed)
                reconnectTimer = setTimeout(connect, 3000)
        }
    }
    connect()
    return () => {
        closed = true
        clearTimeout(reconnectTimer)
        if (ws) { try { ws.close() } catch {} }
    }
}
\end{lstlisting}


\subsection{Реализация клиентского приложения}

Клиентское приложение реализовано на
React~19~\cite{react_official,react_reference_hooks} с
использованием сборщика Vite~8~\cite{vite_official,vite_guide}.
Архитектура построена на функциональных компонентах с
хуками~\cite{react_reference_hooks}: \texttt{useState} для
управления локальным состоянием, \texttt{useEffect} для
побочных эффектов (загрузка данных, WebSocket-подключение),
\texttt{useMemo} для мемоизации вычислений,
\texttt{useRef} для хранения мутабельных ссылок.

Навигация реализована через состояние \texttt{page}
в корневом компоненте \texttt{App.jsx} с персистенцией
в \texttt{localStorage}~\cite{mdn_webstorage}: при
обновлении страницы пользователь остаётся на том же
экране. Компонент \texttt{Shell} формирует макет приложения
(sidebar, topbar, main) с условным отображением
административного пункта меню <<Операторы>> только для
роли \texttt{agronomist}.

В таблице~\ref{tab:pages} приведён перечень страниц
приложения с описанием функциональности.

\begin{longtable}{|l|l|p{6.2cm}|}
    \caption{Страницы клиентского приложения} \label{tab:pages} \\

    \hline
    \textbf{Страница} & \textbf{Роль} & \textbf{Функциональность} \\ \hline
    \endfirsthead

    \tabcont{\ref{tab:pages}}
    \textbf{Страница} & \textbf{Роль} & \textbf{Функциональность} \\ \hline
    \endhead

    \hline
    \endfoot

    Dashboard & все & KPI-карточки со спарклайнами, карта секторов, график активности за 24~часа, лента предупреждений, таблица лидеров (топ-5) \\ \hline
    Sectors   & все & Плитки секторов, панель полива с визуализацией растения, коэффициенты модели и CWSI, обработка от вредителей, телеметрия, назначение оператора, ручная корректировка параметров и инъекция событий (агроном) \\ \hline
    Watering  & все & Полноэкранный пульт полива: визуализация плантации, коэффициенты модели ($K_s$, $K_{s,aer}$, CWSI, $D_r$, GDD) с подсказками-формулами, обработка от вредителей, индикатор поломки оборудования, личный счёт и бейджи оператора \\ \hline
    Plants    & все & Таблица растений с фильтрацией по сектору, добавление и удаление экземпляров \\ \hline
    Reports   & все & Агрегированные KPI по сектору, графики влажности/здоровья и температуры, табличная телеметрия, экспорт XLSX и настройка генератора погоды (агроном) \\ \hline
    Operators & agro & Таблица лидеров операторов с бейджами, ручное начисление баллов и бейджей, группировка секторов по оператору, средние показатели здоровья и эффективности воды \\ \hline
\end{longtable}


\subsection{Реализация геймифицированного интерфейса}

Геймификация реализована через визуальные компоненты,
обеспечивающие наглядную обратную связь о состоянии растений
и создающие мотивацию к своевременному поливу.

\subsubsection{Визуализация растения}

Компонент \texttt{PlantVisual} (листинг~\ref{lst:plant_visual})
представляет собой процедурно генерируемую SVG-иллюстрацию
дерева, параметрически зависящую от текущей влажности почвы
$\theta$ и индекса здоровья $H$:

\begin{itemize}
    \item \textbf{Увядание} --- при $\theta < 25$\% ствол
    и ветви деформируются (наклон пропорционален
    $1 - \theta / 25$), имитируя поникание кроны.
    \item \textbf{Затопление} --- при $\theta > 90$\%
    отображается эллипс водной поверхности у основания
    дерева.
    \item \textbf{Цвет кроны} --- определяется по $H$ через
    цветовое пространство OKLCH: зелёный ($H > 70$),
    жёлто-зелёный ($H > 40$), коричневый ($H \leq 40$).
    \item \textbf{Плоды} --- красные точки отображаются
    только при $H > 55$, визуализируя плодоношение.
    \item \textbf{Полоска влажности} --- горизонтальный
    прогресс-бар в нижней части отображает $\theta / 100$.
\end{itemize}

\begin{lstlisting}[caption={Процедурная визуализация растения (visuals.jsx, фрагмент)}, label={lst:plant_visual}]
export function PlantVisual({ moisture, health,
    watering }) {
    const m = moisture ?? 50
    const h = health ?? 50
    const wilt = m < 25 ? 1 - (m / 25) : 0
    const flood = m > 90 ? (m - 90) / 10 : 0
    const vigor = Math.max(0.3, h / 100)
    const leafColor = h > 70
        ? 'oklch(0.55 0.14 140)'
        : h > 40
            ? 'oklch(0.55 0.14 100)'
            : 'oklch(0.48 0.12 60)'
    return (
        <svg viewBox="0 0 200 200">
          {/* sky, soil, trunk, canopy, fruits */}
          {watering && /* rain drops animation */}
        </svg>
    )
}
\end{lstlisting}

При активации полива отображается анимация падающих капель
с использованием SVG-элемента \texttt{<animate>},
создающая визуальную обратную связь для оператора.

\subsubsection{Кольцевой индикатор здоровья}

Компонент \texttt{HealthRing} реализует кольцевую диаграмму
на базе SVG-элемента \texttt{<circle>} с вычисляемым
атрибутом \texttt{stroke-dashoffset}. Цвет кольца
динамически изменяется: зелёный ($H > 70$), жёлтый
($H > 40$), красный ($H \leq 40$), обеспечивая мгновенную
визуальную индикацию критических состояний.

\subsubsection{Графики телеметрии}

Компоненты \texttt{Sparkline} и \texttt{LineChart} реализуют
визуализацию временных рядов на чистом SVG без внешних
библиотек. \texttt{Sparkline} используется в KPI-карточках
дашборда для отображения тренда последних 48~точек.
\texttt{LineChart} --- полноразмерный график с осями,
сеткой и легендой, адаптивно масштабируемый через
\texttt{ResizeObserver}~\cite{mdn_resizeobserver}.

\subsubsection{Система уведомлений и обратная связь}

При получении WebSocket-события \texttt{notification}
система отображает всплывающее уведомление (\emph{toast})
с текстом сообщения и типом события (засуха, переувлажнение,
критическое здоровье). Уведомления автоматически скрываются
через 4~секунды. Боковая панель <<Центр событий>> отображает
историю уведомлений с возможностью пометить как прочитанное
(индивидуально или все сразу).

\subsubsection{Прогресс-бары и лимиты}

Суточный лимит воды визуализирован прогресс-баром
с цветовой индикацией: зелёный при потреблении менее 85\%
лимита, оранжевый --- при превышении порога. При попытке
полива сверх лимита API возвращает ошибку, а клиент
отображает соответствующее уведомление.

\subsubsection{Система очков, бейджей и таблица лидеров}

Ядром геймификации служит система обучающей аналитики,
начисляющая оператору очки на каждом тике движка (подсистема
\texttt{ScoringTick}). Положительные очки начисляются за
поддержание индекса здоровья $H > 0{,}7$ и эффективный полив
без стресса ($K_s > 0{,}9$, $K_{s,aer} > 0{,}9$); штрафы ---
за допущенную засуху ($K_s < 0{,}3$) и переувлажнение
($K_{s,aer} < 0{,}5$). За устойчивые серии присуждаются
бейджи: <<Хранитель воды>> (7~суток без перелива),
<<Зелёный мастер>> ($H > 0{,}9$ в течение 14~суток) и
<<Кризис-менеджер>> (преодоление трёх случайных событий
подряд без потери здоровья). Серии отслеживаются счётчиками
на секторе (\texttt{healthy\_streak}, \texttt{safe\_streak},
\texttt{crisis\_streak}), а накопленные результаты
агрегируются в таблице \texttt{training\_scores} (суммарный
балл, бейджи в формате JSONB, среднее здоровье, эффективность
использования воды).

На клиенте личный счёт и полученные бейджи отображаются на
пульте полива, а агроному на странице <<Операторы>> доступна
таблица лидеров (\texttt{GET /api/leaderboard}), ранжирующая
операторов по суммарному баллу с указанием среднего здоровья,
эффективности воды и заслуженных бейджей. Сокращённая таблица
лидеров (топ-5) дублируется на дашборде.

Помимо автоматического начисления, реализовано
\emph{ручное поощрение наставником}
(\texttt{POST /api/training/\{userId\}/award}): в карточке
каждого оператора на странице <<Операторы>> агроном может
начислить или списать баллы (значение может быть
отрицательным) и выдать бейджи. Бейджи объединяются с уже
имеющимися без дублирования; при отсутствии у оператора
записи аналитики она создаётся. Маршрут защищён ролью
\texttt{agronomist} и проверяет, что цель начисления имеет
роль \texttt{operator}.


\subsection{Реализация управления поливом}

Управление поливом реализовано как взаимодействие серверного
обработчика \texttt{WateringHandler} и клиентского интерфейса
полива, представленного как компактной панелью на странице
<<Секторы>>, так и отдельным полноэкранным пультом полива
(страница <<Полив>>).

\subsubsection{Серверная обработка полива}

При получении запроса \texttt{POST /api/water} обработчик
выполняет следующую последовательность:
\begin{enumerate}
    \item Извлечение \texttt{user\_id} и \texttt{role} из
    контекста JWT.
    \item Проверка актуальности сессии: если пользователь из
    токена отсутствует в БД (устаревший JWT, например после
    пересоздания базы), возвращается HTTP~401, и клиент
    инициирует повторный вход.
    \item Валидация входных данных: \texttt{sector\_id} и
    \texttt{volume\_liters} (от 0 до 1000~литров).
    \item Проверка назначения: оператор может поливать
    только закреплённый за ним сектор.
    \item Проверка доступности оборудования: при активной
    поломке (\texttt{equipment\_locked\_ticks} $> 0$) полив
    отклоняется с кодом HTTP~409 (Conflict).
    \item Проверка суточного лимита: если объём
    \(\text{water\_consumed} + V\) превышает
    \(\text{daily\_water\_limit}\), возвращается ошибка.
    \item Запись в журнал полива (\texttt{watering\_logs}).
    \item Пересчёт дефицита влаги: $D_r = D_r - V$, где
    объём полива $V$ (в литрах) учитывается как эквивалентный
    слой воды в миллиметрах (1~л~$\approx$~1~мм) для баланса
    игровой механики. Если объём превышает текущий дефицит,
    излишек поднимает влажность почвы выше полевой влагоёмкости,
    моделируя переувлажнение.
    \item Обновление \texttt{water\_consumed},
    \texttt{last\_watered\_at} и сохранение в БД.
    \item Широковещательная рассылка события
    \texttt{sector:watered} через WebSocket.
\end{enumerate}

Отдельным действием оператора является \emph{обработка от
вредителей} (\texttt{POST /api/sectors/\{id\}/treat}): при
активном заражении (\texttt{pest\_active}) обработчик снимает
его, возвращая сектор в нормальное состояние. Действие
доступно назначенному оператору и агроному и проверяет права
доступа аналогично поливу.

\subsubsection{Клиентский интерфейс полива}

Панель полива интегрирована в страницу <<Секторы>> и
включает:
\begin{itemize}
    \item Кнопки быстрого выбора объёма (5, 10, 20, 40~л)
    и ползунок \texttt{<input type="range">} для точной
    настройки от 1 до 100~л.
    \item Кнопку <<Полить сектор>> с анимацией нажатия
    (\texttt{transform: scale(0.98)}).
    \item Визуализацию \texttt{PlantVisual} с динамическим
    отображением состояния растения.
    \item Прогресс-бары влажности почвы и суточного лимита.
    \item Блокировку кнопки, если сектор закреплён за другим
    оператором (либо при поломке оборудования), с
    пояснительным текстом.
\end{itemize}

Полноэкранный пульт полива (страница <<Полив>>) дополнительно
отображает коэффициенты математической модели ($K_s$,
$K_{s,aer}$, CWSI, $D_r$, GDD) со всплывающими
подсказками-формулами, индикатор поломки оборудования,
баннер нашествия вредителей с кнопкой <<Обработать>> и
личный счёт оператора с заработанными бейджами.


\subsection{Реализация импорта и экспорта данных}

В соответствии с задачей~6 (подраздел~1.3) реализованы
механизмы импорта и экспорта конфигурации секторов, а также
выгрузки телеметрии для офлайн-анализа.

\subsubsection{Экспорт в CSV}

Обработчик \texttt{ExportSectors}
(листинг~\ref{lst:export_csv}) формирует CSV-файл с
заголовками \texttt{Content-Type: text/csv; charset=utf-8} и
\texttt{Content-Disposition: attachment} для автоматической
загрузки браузером. Файл начинается с метки порядка байтов
UTF-8 (BOM) и использует разделитель <<точка с запятой>>,
что обеспечивает корректное отображение кириллических
названий и разбиение по столбцам в табличных процессорах
с русской локалью. Каждая строка содержит 12~полей сектора:
название, площадь, влажность, температуру, индекс здоровья,
кумулятивные GDD, фенофазу, коэффициенты $K_s$ и $K_{s,aer}$,
дефицит $D_r$, суточный лимит воды и статус.

\begin{lstlisting}[caption={Экспорт секторов в CSV (handler/files.go)}, label={lst:export_csv}]
func (h *FileHandler) ExportSectors(
    w http.ResponseWriter, r *http.Request) {
    sectors, err := h.store.ListSectors(r.Context())
    if err != nil {
        writeJSON(w, http.StatusInternalServerError,
            map[string]string{"error": err.Error()})
        return
    }
    w.Header().Set("Content-Type",
        "text/csv; charset=utf-8")
    w.Header().Set("Content-Disposition",
        "attachment; filename=sectors.csv")
    w.Write(utf8BOM) // UTF-8 BOM for Excel
    writer := csv.NewWriter(w)
    writer.Comma = ';'
    defer writer.Flush()
    writer.Write([]string{"name", "area_sqm",
        "soil_moisture", "temperature",
        "health_index", "gdd_cumulative",
        "phenophase", "ks_water", "ks_aeration",
        "deficit_dr", "daily_water_limit", "status"})
    for _, s := range sectors {
        writer.Write([]string{
            s.Name,
            fmt.Sprintf("%.2f", s.AreaSqm),
            fmt.Sprintf("%.2f", s.SoilMoisture),
            /* ... remaining fields ... */
        })
    }
}
\end{lstlisting}

\subsubsection{Импорт из CSV}

Импорт осуществляется через \texttt{multipart/form-data}
(\texttt{POST /api/import/sectors}) и устойчив к различным
форматам входного файла: срезается метка UTF-8 BOM,
автоматически определяется разделитель (<<точка с запятой>>
или запятая), допускается десятичная запятая. Обработчик
распознаёт строку заголовка и сопоставляет столбцы по именам
(в произвольном порядке и в любом подмножестве), а при
отсутствии заголовка использует позиционное соответствие
<<название, площадь>>. Отсутствующие параметры симуляции
заполняются корректными значениями по умолчанию (влажность
60\%, индекс здоровья 1{,}0, $K_s = K_{s,aer} = 1{,}0$,
суточный лимит 500~л), а присутствующие --- восстанавливаются,
после чего производные коэффициенты ($K_s$, $K_{s,aer}$,
влажность~$\leftrightarrow$~дефицит, фенофаза по GDD)
пересчитываются для внутренней согласованности сектора.
В ответе возвращается количество импортированных и
пропущенных записей.

На клиенте импорт реализован через скрытый элемент
\texttt{<input type="file" accept=".csv">}, активируемый
кнопкой <<Импорт CSV>> (доступна только агроному); выбранный
файл передаётся в обработчик напрямую и загружается через
\texttt{FormData API} без перезагрузки страницы.

\subsubsection{Экспорт телеметрии в XLSX}

Помимо CSV, реализован экспорт временного ряда телеметрии
выбранного сектора в формат XLSX
(\texttt{GET /api/export/telemetry/\{id\}}), удобный для
последующего анализа агрономом в табличном процессоре.
Генерация файла выполнена в собственном пакете
\texttt{internal/xlsx}, написанном исключительно на
стандартной библиотеке Go (\texttt{archive/zip} и
\texttt{encoding/xml}) без сторонних зависимостей:
формируется минимальный валидный документ OOXML
SpreadsheetML с одним листом, содержащим строковые и
числовые ячейки (время, влажность почвы, температура и
индекс здоровья). На клиенте выгрузка инициируется кнопкой
<<XLSX>> на странице <<Секторы>> и доступна только агроному.


\subsection{Генерация начальных данных и тестирование}

\subsubsection{Автоматическая генерация тестовых данных}

Для демонстрации и тестирования приложения реализована
функция \texttt{Seed} (пакет \texttt{internal/storage}),
автоматически заполняющая пустую базу данных при первом
запуске. Функция создаёт:
\begin{itemize}
    \item 1~агронома и 3~операторов с предустановленными
    учётными данными;
    \item 6~секторов плантации с рандомизированными
    начальными параметрами (площадь 500--2500~м$^2$,
    влажность 45--70\%, температура 24--30~°C, индекс
    здоровья 0{,}7--1{,}0) и автоматическим назначением
    операторов;
    \item от 18 до 42~растений (\textit{Litchi chinensis},
    в том числе разновидность \textit{var.~philippinensis})
    со случайным распределением по секторам (3--7 на сектор);
    \item историю полива за 7~дней (2--5~событий в сутки
    на сектор);
    \item 144~точки телеметрии на сектор (24~часа с
    интервалом 10~минут) с реалистичным случайным дрейфом;
    \item 6--18~уведомлений различных типов.
\end{itemize}

\subsubsection{Функциональное тестирование}

Тестирование проводилось в контейнерном окружении
(\texttt{docker compose up}) по следующим сценариям:

\begin{enumerate}
    \item \textbf{Аутентификация и авторизация}: вход под
    обеими ролями, попытка доступа к административным
    маршрутам от имени оператора (ожидание HTTP~403),
    регистрация нового пользователя.
    \item \textbf{Работа движка симуляции}: наблюдение
    за изменением влажности, температуры и индекса здоровья
    в реальном времени; верификация формул через сопоставление
    значений $K_s$, $K_{s,aer}$ и $D_r$ с ожидаемыми
    (ручной расчёт по формулам подраздела~2.3).
    \item \textbf{Управление поливом}: полив сектора с
    различными объёмами, проверка суточного лимита,
    попытка полива чужого сектора (ожидание HTTP~403),
    визуальная проверка анимации и обновления прогресс-баров.
    \item \textbf{WebSocket}: подключение нескольких
    клиентов одновременно, верификация синхронного
    обновления состояния секторов на всех экранах.
    \item \textbf{Уведомления}: проверка генерации
    уведомлений при достижении пороговых значений $K_s$
    и $H$, доставка через WebSocket, пометка как прочитано.
    \item \textbf{Импорт/экспорт}: экспорт секторов в CSV,
    импорт из файла, экспорт телеметрии в XLSX и открытие
    файла в табличном процессоре, проверка корректности данных.
    \item \textbf{Случайные события}: инъекция нашествия
    вредителей и поломки оборудования, проверка постоянного
    урона здоровью при заражении, блокировки полива при
    поломке, выполнение обработки сектора и снятие заражения.
    \item \textbf{Геймификация}: верификация начисления очков
    и присуждения бейджей при поддержании здоровья и
    эффективного полива, штрафов за засуху и переувлажнение,
    корректность ранжирования операторов в таблице лидеров,
    ручное начисление баллов и бейджей агрономом.
\end{enumerate}

В ходе тестирования подтверждена корректность реализации
математической модели: при отсутствии полива дефицит влаги
$D_r$ монотонно возрастает за счёт эвапотранспирации,
коэффициент $K_s$ уменьшается, индекс здоровья деградирует;
своевременный полив снижает $D_r$, восстанавливает $K_s$
и активирует восстановление $H$ через коэффициент
$\beta_{rec}$. Переходы между состояниями автомата
(нормальное $\to$ засуха $\to$ критическое $\to$
восстановление) воспроизводились стабильно.
